% GENERATED FILE. Edit canonical lessons, then run npm run generate:books.
% canonical-source-hash: fnv1a64:fd22cef00b404377
% canonical-lessons: ML-C160-koni, ML-C160-varanta, ML-C160-malavillu, ML-C160-vellaccattam, ML-C160-kallu

\chapter{Ladder, Verandah, Rainbow, Waterfall, Stone}
\label{ch:ml-ladder-verandah-rainbow-waterfall-stone}
\hlchaptermodality{160}

\begin{chapteropening}
\textbf{By the end of this chapter:} \emph{I can say a ladder, a verandah, a rainbow, a waterfall and a stone.}

Complete the last lesson of chapter 160: I can say a ladder, a verandah, a rainbow, a waterfall and a stone.
\end{chapteropening}

\section[\texorpdfstring{\ml{കോണി} (kōṇi)}{കോണി (kōṇi)}]{\ml{കോണി} (kōṇi) — a ladder}

\label{lesson:ML-C160-koni}



\begin{quote}
Before the new one: say the Malayalam for a fence, then the Malayalam for a gate.

Say the four seasons, then the one that shapes Kerala most. (\emph{Maḻakkālaṁ}, the rains.)
\end{quote}

\subsection*{You'll want to know: \ml{കോണി}}
\textbf{\ml{കോണി}} — \emph{kōṇi} — \textquotedblleft{}a ladder\textquotedblright{}.

\begin{cousinweb}
\textbf{\ml{കോണി}} (\emph{kōṇi}) is of uncertain origin. DEDR lists a Malayalam \emph{kōṇi}, but meaning \textquotedblleft{}gunny bag\textquotedblright{} (DEDR appendix 33), a different word.
\end{cousinweb}

\begin{grammarlens}[title={What you've built}]
One more word for this chapter.
\end{grammarlens}

\subsection*{Your turn}
\begin{itemize}
\raggedright
  \item \emph{Say it:} \emph{kōṇi}
  \item \emph{Say it:} \emph{kōṇi}, once more
  \item \emph{Say it:} say \emph{kōṇi}
  \item \emph{Recall:} say the Malayalam for a fence, then the Malayalam for a gate, then say \emph{kōṇi} again
\end{itemize}

\begin{tcolorbox}[breakable,colback=teal!4,colframe=teal!35!black,title={Before you move on}]
What does \ml{കോണി} mean? (\textquotedblleft{}A ladder\textquotedblright{}.) Say it once more.
\end{tcolorbox}
\section[\texorpdfstring{\ml{വരാന്ത} (varānta)}{വരാന്ത (varānta)}]{\ml{വരാന്ത} (varānta) — a verandah}

\label{lesson:ML-C160-varanta}



\begin{quote}
Before the new one: say the Malayalam for a gate, then the Malayalam for a ladder.
\end{quote}

\subsection*{You'll want to know: \ml{വരാന്ത}}
\textbf{\ml{വരാന്ത}} — \emph{varānta} — \textquotedblleft{}a verandah\textquotedblright{}.

\begin{cousinweb}
\textbf{\ml{വരാന്ത}} (\emph{varānta}) $\leftarrow$ Portuguese \emph{varanda}, \textquotedblleft{}balcony, verandah\textquotedblright{}; English \emph{verandah} came by the same road.
\end{cousinweb}

\begin{grammarlens}[title={What you've built}]
One more word for this chapter.
\end{grammarlens}

\subsection*{Your turn}
\begin{itemize}
\raggedright
  \item \emph{Say it:} \emph{varānta}
  \item \emph{Say it:} \emph{varānta}, once more
  \item \emph{Say it:} say \emph{varānta}
  \item \emph{Recall:} say the Malayalam for a gate, then the Malayalam for a ladder, then say \emph{varānta} again
\end{itemize}

\begin{tcolorbox}[breakable,colback=teal!4,colframe=teal!35!black,title={Before you move on}]
What does \ml{വരാന്ത} mean? (\textquotedblleft{}A verandah\textquotedblright{}.) Say it once more.
\end{tcolorbox}
\section[\texorpdfstring{\ml{മഴവില്ല്} (maḻavillŭ)}{മഴവില്ല് (maḻavillŭ)}]{\ml{മഴവില്ല്} (maḻavillŭ) — a rainbow}

\label{lesson:ML-C160-malavillu}



\begin{quote}
Before the new one: say the Malayalam for a ladder, then the Malayalam for a verandah.
\end{quote}

\subsection*{You'll want to know: \ml{മഴവില്ല്}}
\textbf{\ml{മഴവില്ല്}} — \emph{maḻavillŭ} — \textquotedblleft{}a rainbow\textquotedblright{}.

\begin{cousinweb}
\textbf{\ml{മഴവില്ല്}} (\emph{maḻavillŭ}) is native Dravidian, \textquotedblleft{}the rain bow\textquotedblright{}, built exactly like the English: \emph{maḻa}, \textquotedblleft{}rain\textquotedblright{} (DEDR 4753, with Tamil \emph{maḻai}), and \emph{villŭ}, \textquotedblleft{}bow\textquotedblright{} (DEDR 5422, with Tamil \emph{vil}).
\end{cousinweb}

\begin{grammarlens}[title={What you've built}]
One more word for this chapter.
\end{grammarlens}

\subsection*{Your turn}
\begin{itemize}
\raggedright
  \item \emph{Say it:} \emph{maḻavillŭ}
  \item \emph{Say it:} \emph{maḻavillŭ}, once more
  \item \emph{Say it:} say \emph{maḻavillŭ}
  \item \emph{Recall:} say the Malayalam for a ladder, then the Malayalam for a verandah, then say \emph{maḻavillŭ} again
\end{itemize}

\begin{tcolorbox}[breakable,colback=teal!4,colframe=teal!35!black,title={Before you move on}]
What does \ml{മഴവില്ല്} mean? (\textquotedblleft{}A rainbow\textquotedblright{}.) Say it once more.
\end{tcolorbox}
\section[\texorpdfstring{\ml{വെള്ളച്ചാട്ടം} (veḷḷaccāṭṭaṁ)}{വെള്ളച്ചാട്ടം (veḷḷaccāṭṭaṁ)}]{\ml{വെള്ളച്ചാട്ടം} (veḷḷaccāṭṭaṁ) — a waterfall}

\label{lesson:ML-C160-vellaccattam}



\begin{quote}
Before the new one: say the Malayalam for a verandah, then the Malayalam for a rainbow.

Say the first three Malayalam months. (\emph{Ciṅṅaṁ}, \emph{kanni}, \emph{tulāṁ}.)
\end{quote}

\subsection*{You'll want to know: \ml{വെള്ളച്ചാട്ടം}}
\textbf{\ml{വെള്ളച്ചാട്ടം}} — \emph{veḷḷaccāṭṭaṁ} — \textquotedblleft{}a waterfall\textquotedblright{}.

\begin{cousinweb}
\textbf{\ml{വെള്ളച്ചാട്ടം}} (\emph{veḷḷaccāṭṭaṁ}) is native Dravidian, \textquotedblleft{}the water's leap\textquotedblright{}: \emph{veḷḷam}, \textquotedblleft{}water\textquotedblright{} (DEDR 5503), and \emph{cāṭṭam}, \textquotedblleft{}a leap\textquotedblright{} (DEDR 2440, with Tamil \emph{cāṭṭam}).
\end{cousinweb}

\begin{grammarlens}[title={What you've built}]
One more word for this chapter.
\end{grammarlens}

\subsection*{Your turn}
\begin{itemize}
\raggedright
  \item \emph{Say it:} \emph{veḷḷaccāṭṭaṁ}
  \item \emph{Say it:} \emph{veḷḷaccāṭṭaṁ}, once more
  \item \emph{Say it:} say \emph{veḷḷaccāṭṭaṁ}
  \item \emph{Recall:} say the Malayalam for a verandah, then the Malayalam for a rainbow, then say \emph{veḷḷaccāṭṭaṁ} again
\end{itemize}

\begin{tcolorbox}[breakable,colback=teal!4,colframe=teal!35!black,title={Before you move on}]
What does \ml{വെള്ളച്ചാട്ടം} mean? (\textquotedblleft{}A waterfall\textquotedblright{}.) Say it once more.
\end{tcolorbox}
\section[\texorpdfstring{\ml{കല്ല്} (kallŭ)}{കല്ല് (kallŭ)}]{\ml{കല്ല്} (kallŭ) — a stone}

\label{lesson:ML-C160-kallu}



\begin{quote}
Before the new one: say the Malayalam for a rainbow, then the Malayalam for a waterfall.
\end{quote}

\subsection*{You'll want to know: \ml{കല്ല്}}
\textbf{\ml{കല്ല്}} — \emph{kallŭ} — \textquotedblleft{}a stone\textquotedblright{}.

\begin{cousinweb}
\textbf{\ml{കല്ല്}} (\emph{kallŭ}) is native Dravidian, one word with Tamil \emph{kal}, Kannada \emph{kal} and Telugu \emph{kallu} (DEDR 1298).
\end{cousinweb}

\begin{grammarlens}[title={What you've built}]
One more word for this chapter.
\end{grammarlens}

\subsection*{Your turn}
\begin{itemize}
\raggedright
  \item \emph{Say it:} \emph{kallŭ}
  \item \emph{Say it:} \emph{kallŭ}, once more
  \item \emph{Say it:} say \emph{kallŭ}
  \item \emph{Recall:} say the Malayalam for a rainbow, then the Malayalam for a waterfall, then say \emph{kallŭ} again
\end{itemize}

\begin{tcolorbox}[breakable,colback=teal!4,colframe=teal!35!black,title={Before you move on}]
What does \ml{കല്ല്} mean? (\textquotedblleft{}A stone\textquotedblright{}.) Say it once more.
\end{tcolorbox}
